\section{Mapping Musical Primitives to Programming Constructs}\label{sec:musical-mapping}

To bridge the gap between abstract musical theory and formal computer science, the proposed DSL maps fundamental musical entities to specific programming constructs.
This mapping ensures that the language remains expressive for musicians while retaining the type safety and structural clarity of a compiled language.

\subsection{Musical Entities as Data Types}

The language defines several core primitives that serve as the building blocks of a composition:
\begin{itemize}
    \item \textbf{Notes:} Represented as pitch-class and octave pairs (e.g., \texttt{C4}, \texttt{F\#2}). These are treated as constant literals with unique numeric identifiers in the underlying MIDI protocol.
    \item \textbf{Chords:} Modeled as collections of simultaneous notes. In the DSL, a chord is a single unit of duration, where the temporal span is determined by its constituent notes.
    \item \textbf{Sequences:} Represented as ordered lists of musical events (notes, chords, or rests). This structure maps directly to the concept of a musical "bar" or "cell."
    \item \textbf{Rests:} Unlike many languages that use a \texttt{sleep} function to pause execution, this DSL treats a rest as a durational value. A rest is a "silent note" that occupies space in the timeline, maintaining the structural alignment of the composition.
\end{itemize}

\subsection{Compositional Abstractions}

The hierarchical organization of music is reflected in the language's scoping and container constructs:

\subsubsection{Patterns: The Unit of Reusability}
A \texttt{pattern} is the primary abstraction mechanism in the DSL.
Semantically, it is equivalent to a function definition.
It allows a composer to encapsulate a motif and assign it a name.
Patterns can accept parameters, enabling the creation of dynamic musical ideas—such as a drum fill that varies based on a boolean flag or a melodic line that accepts a sequence as a template.

\subsubsection{Tracks: Instrumental Identity}
A \texttt{track} serves as the top-level execution context for a specific instrument.
In the compilation process, each track maps to a unique MIDI channel.
This ensures that independent musical lines (e.g., a piano melody and a bassline) are kept separate and assigned to distinct sound patches in the final artifact.

\subsubsection{Voices: Horizontal Parallelism}
Within a single track, the \texttt{voice} construct allows for polyphonic composition.
While a track represents a single "performer," a voice represents a single "line of counterpoint."
Multiple voices can be defined to start at the same time or at different offsets within the same track, allowing for complex harmonic structures without the need for manual delta-time calculations.
Voices enable the composer to think in terms of horizontal layers rather than just vertical chords.
